\documentclass[12pt]{article}
\usepackage[utf8]{inputenc}
\usepackage[T1]{fontenc}
\usepackage{hyperref}
\title{Liana-Seti Language Resources for K-12 Curriculum Development}
\author{Compiled by Claude Sonnet 5.5 (claude-sonnet-5-5) for Andr\'as Kornai}
\date{}
\begin{document}
\maketitle

\section{Introduction}
Liana-Seti (ISO 639-3: ste; Glottocode lian1255) is an Austronesian language of Seram island in Maluku, Indonesia. Wikipedia places it in Central Malayo-Polynesian, under East Seram. It is also listed under the names Kobi and Uhei Kachlakan, which it shares with the neighbouring Benggoi language. Speaker figures vary by source. One figure is about 3,000 speakers in eight villages, and Joshua Project gives a population of 4,900 for the Liana people. Little has been published on the language.

\section{Primary Educational Resources}
No primer, textbook or school programme for Liana-Seti was found in the search results.

\section{Linguistic and Reference Materials}
\begin{itemize}
\item Wikipedia, ``Liana language'': \url{https://en.wikipedia.org/wiki/Liana_language}
\item Glottolog, Seti languoid (classification and bibliography): \url{https://glottolog.org/resource/languoid/id/seti1246.bigmap.html}
\item Numeral systems database, Liana-Seti (a decimal counting system with example number words): \url{https://lingweb.eva.mpg.de/channumerals/Liana-Seti.htm}
\end{itemize}

\section{Religious and Community Resources}
\begin{itemize}
\item Joshua Project, Liana in Indonesia (a general people-group profile; it is not well sourced, so treat it with caution): \url{https://joshuaproject.net/people_groups/13052}
\end{itemize}
No scripture translation or community language organisation was found.

\section{Digital and Audio Resources}
No audio, video or corpus resources were found.

\section{Curriculum Development Context}
The language is endangered, but it is still used by some people in every generation. The search results show no language documentation beyond basic reference data. Curriculum work would have to begin with community-led documentation. The surrounding language environment (Indonesian and Malay-based regional varieties) will shape any bilingual design.
\end{document}